\documentclass[12pt]{article}
\usepackage[utf8]{inputenc}
\usepackage[T1]{fontenc}
\usepackage{hyperref}
\title{Situ Language Resources for K-12 Curriculum Development}
\author{Compiled by Claude Sonnet 5 (claude-sonnet-5) for Andr\'as Kornai}
\date{}
\begin{document}
\maketitle
\section{Introduction}
Situ (ISO 639-3 tzi) is a Gyalrongic (Rgyalrong) language of Sichuan, China. Search found only academic and encyclopedic material. It found no teaching materials.
\section{Primary Educational Resources}
None found.
\section{Linguistic and Reference Materials}
\begin{itemize}
\item Situ language overview (Wikipedia). \url{https://en.wikipedia.org/wiki/Situ_language}
\item Gyalrong languages overview. \url{https://en.wikipedia.org/wiki/Gyalrong_languages}
\item Gyalrongic languages overview. \url{https://en.wikipedia.org/wiki/Gyalrongic_languages}
\item Zhang Shuya and Fan Jingming, ``Brag-bar kinship system in synchronic and diachronic perspectives'', BSOAS. It covers the Brag-bar dialect of Situ. \url{https://cambridge.org/core/journals/bulletin-of-the-school-of-oriental-and-african-studies/article/abs/bragbar-kinship-system-in-synchronic-and-diachronic-perspectives/BB56B51BD5591C33215224864FFA0CFE}
\item A Rgyalrongic lexical database article. \url{https://pmc.ncbi.nlm.nih.gov/articles/PMC10446046}
\end{itemize}
\section{Religious and Community Resources}
None found.
\section{Digital and Audio Resources}
No audio or software resources were found. The lexical database article above may be of use.
\section{Curriculum Development Context}
No school programme or standardized orthography for Situ turned up in search. Gates's \emph{Situ in Situ} (Lincom, 2014) is cited by secondary sources, but I did not see a direct URL for it. A curriculum would need new, community-led materials.
\end{document}